\documentclass[10pt,a4paper]{amsart}
\usepackage{amsmath,amssymb,amsthm,graphicx}
\usepackage{verbatim}
\newenvironment{reference}{\comment}{\endcomment}
%\newenvironment{reference}{}{}
\newenvironment{slogan}{\comment}{\endcomment}
\newenvironment{history}{\comment}{\endcomment}

% For commutative diagrams we use Xy-pic
\usepackage[all]{xy}

% We use 2cell for 2-commutative diagrams.
\xyoption{2cell}
% Theorem environments.
%
\theoremstyle{plain}
\newtheorem{theorem}[subsection]{Theorem}
\newtheorem{proposition}[subsection]{Proposition}
\newtheorem{lemma}[subsection]{Lemma}

\theoremstyle{definition}
\newtheorem{definition}[subsection]{Definition}
\newtheorem{example}[subsection]{Example}
\newtheorem{exercise}[subsection]{Exercise}
\newtheorem{situation}[subsection]{Situation}

\theoremstyle{remark}
\newtheorem{remark}[subsection]{Remark}
\newtheorem{remarks}[subsection]{Remarks}

\numberwithin{equation}{subsection}

% Macros
%
\def\lim{\mathop{\mathrm{lim}}\nolimits}
\def\colim{\mathop{\mathrm{colim}}\nolimits}
\def\Spec{\mathop{\mathrm{Spec}}}
\def\Hom{\mathop{\mathrm{Hom}}\nolimits}
\def\Ext{\mathop{\mathrm{Ext}}\nolimits}
\def\SheafHom{\mathop{\mathcal{H}\!\mathit{om}}\nolimits}
\def\SheafExt{\mathop{\mathcal{E}\!\mathit{xt}}\nolimits}
\def\Sch{\mathit{Sch}}
\def\Mor{\mathop{\mathrm{Mor}}\nolimits}
\def\Ob{\mathop{\mathrm{Ob}}\nolimits}
\def\Sh{\mathop{\mathit{Sh}}\nolimits}
\def\NL{\mathop{N\!L}\nolimits}
\def\CH{\mathop{\mathrm{CH}}\nolimits}
\def\proetale{{pro\text{-}\acute{e}tale}}
\def\etale{{\acute{e}tale}}
\def\QCoh{\mathit{QCoh}}
\def\Ker{\mathop{\mathrm{Ker}}}
\def\Im{\mathop{\mathrm{Im}}}
\def\Coker{\mathop{\mathrm{Coker}}}
\def\Coim{\mathop{\mathrm{Coim}}}

% Boxtimes
%
\DeclareMathSymbol{\boxtimes}{\mathbin}{AMSa}{"02}

%
% Macros for moduli stacks/spaces
%
\def\QCohstack{\mathcal{QC}\!\mathit{oh}}
\def\Cohstack{\mathcal{C}\!\mathit{oh}}
\def\Spacesstack{\mathcal{S}\!\mathit{paces}}
\def\Quotfunctor{\mathrm{Quot}}
\def\Hilbfunctor{\mathrm{Hilb}}
\def\Curvesstack{\mathcal{C}\!\mathit{urves}}
\def\Polarizedstack{\mathcal{P}\!\mathit{olarized}}
\def\Complexesstack{\mathcal{C}\!\mathit{omplexes}}
% \Pic is the operator that assigns to X its picard group, usage \Pic(X)
% \Picardstack_{X/B} denotes the Picard stack of X over B
% \Picardfunctor_{X/B} denotes the Picard functor of X over B
\def\Pic{\mathop{\mathrm{Pic}}\nolimits}
\def\Picardstack{\mathcal{P}\!\mathit{ic}}
\def\Picardfunctor{\mathrm{Pic}}
\def\Deformationcategory{\mathcal{D}\!\mathit{ef}}
\usepackage{cleveref}
\crefname{originalreference}{original reference}{original references}
\Crefname{originalreference}{Original reference}{Original references}
\makeatletter
\newlabel{spaces-divisors-definition-relatively-ample}{{0D31}{source}{}{sourcepointer-1}{}}
\newlabel{spaces-divisors-definition-relatively-ample@cref}{{[originalreference][0][]0D31}{[1][0][]source}{}{}{}}
\makeatother
\crefname{originalreference}{original reference}{original references}
\Crefname{originalreference}{Original reference}{Original references}
\makeatletter
\newlabel{spaces-divisors-lemma-relatively-ample-local}{{0D36}{source}{}{sourcepointer-2}{}}
\newlabel{spaces-divisors-lemma-relatively-ample-local@cref}{{[originalreference][0][]0D36}{[1][0][]source}{}{}{}}
\makeatother
\crefname{originallemma}{original lemma}{original lemmas}
\Crefname{originallemma}{Original lemma}{Original lemmas}
\makeatletter
\newlabel{lemma-descending-property-ample}{{0D3C}{source}{}{sourcepointer-3}{}}
\newlabel{lemma-descending-property-ample@cref}{{[originallemma][0][]0D3C}{[1][0][]source}{}{}{}}
\makeatother
\crefname{originalreference}{original reference}{original references}
\Crefname{originalreference}{Original reference}{Original references}
\makeatletter
\newlabel{spaces-limits-lemma-proper-limit-of-proper-finite-presentation-noetherian}{{0A0X}{source}{}{sourcepointer-4}{}}
\newlabel{spaces-limits-lemma-proper-limit-of-proper-finite-presentation-noetherian@cref}{{[originalreference][0][]0A0X}{[1][0][]source}{}{}{}}
\makeatother
\crefname{originalreference}{original reference}{original references}
\Crefname{originalreference}{Original reference}{Original references}
\makeatletter
\newlabel{spaces-limits-lemma-limit-is-affine}{{07SQ}{source}{}{sourcepointer-5}{}}
\newlabel{spaces-limits-lemma-limit-is-affine@cref}{{[originalreference][0][]07SQ}{[1][0][]source}{}{}{}}
\makeatother
\crefname{originalreference}{original reference}{original references}
\Crefname{originalreference}{Original reference}{Original references}
\makeatletter
\newlabel{spaces-limits-lemma-descend-invertible-modules}{{0D2X}{source}{}{sourcepointer-6}{}}
\newlabel{spaces-limits-lemma-descend-invertible-modules@cref}{{[originalreference][0][]0D2X}{[1][0][]source}{}{}{}}
\makeatother
\crefname{originalreference}{original reference}{original references}
\Crefname{originalreference}{Original reference}{Original references}
\makeatletter
\newlabel{spaces-limits-lemma-limit-is-scheme}{{07SR}{source}{}{sourcepointer-7}{}}
\newlabel{spaces-limits-lemma-limit-is-scheme@cref}{{[originalreference][0][]07SR}{[1][0][]source}{}{}{}}
\makeatother
\crefname{originalreference}{original reference}{original references}
\Crefname{originalreference}{Original reference}{Original references}
\makeatletter
\newlabel{limits-lemma-limit-ample}{{09MT}{source}{}{sourcepointer-8}{}}
\newlabel{limits-lemma-limit-ample@cref}{{[originalreference][0][]09MT}{[1][0][]source}{}{}{}}
\makeatother
\crefname{originalreference}{original reference}{original references}
\Crefname{originalreference}{Original reference}{Original references}
\makeatletter
\newlabel{spaces-divisors-lemma-ample-in-neighbourhood}{{0D3A}{source}{}{sourcepointer-9}{}}
\newlabel{spaces-divisors-lemma-ample-in-neighbourhood@cref}{{[originalreference][0][]0D3A}{[1][0][]source}{}{}{}}
\makeatother
\crefname{originalreference}{original reference}{original references}
\Crefname{originalreference}{Original reference}{Original references}
\makeatletter
\newlabel{morphisms-lemma-ample-base-change}{{0893}{source}{}{sourcepointer-10}{}}
\newlabel{morphisms-lemma-ample-base-change@cref}{{[originalreference][0][]0893}{[1][0][]source}{}{}{}}
\makeatother
\crefname{originalreference}{original reference}{original references}
\Crefname{originalreference}{Original reference}{Original references}
\makeatletter
\newlabel{morphisms-lemma-pullback-ample-tensor-relatively-ample}{{0892}{source}{}{sourcepointer-11}{}}
\newlabel{morphisms-lemma-pullback-ample-tensor-relatively-ample@cref}{{[originalreference][0][]0892}{[1][0][]source}{}{}{}}
\makeatother
\title{Ampleness spreads out from fibres along a proper morphism of algebraic spaces}
\author{The Stacks Project Authors}
\date{stacks-project:spaces-descent.tex:local-snapshot-20261012; GNU FDL 1.2 or later}
\begin{document}
\maketitle
\noindent\textit{Editorial note.} All delivered mathematical text is the author's own, copied verbatim from the pinned Stacks Project chapter file spaces-descent.tex; no statement or proof step was rewritten, reordered or removed.\par\medskip
\noindent\textit{Editorial note.} The author closes the final step with three named ingredients ('hints'); the composition is the affine-pullback transfer of ampleness from \$X\_i\$ to \$X\$. Flagged for the reviewer: this closing step is hint-form rather than fully written out, while everything before it is explicit.\par\medskip
\noindent\textit{Allowed standard prerequisites (editorial).}\begin{itemize}
\item fpqc locality of relative ampleness. Stacks Project, Descent on Spaces, Lemma lemma-descending-property-ample. Reduces to residue fields of points.
\item Limits of algebraic spaces. Stacks Project, Limits of Spaces, Lemmas spaces-limits-lemma-proper-limit-of-proper-finite-presentation-noetherian, spaces-limits-lemma-limit-is-affine, spaces-limits-lemma-descend-invertible-modules, spaces-limits-lemma-limit-is-scheme. Supplies the approximation system and its shrinking steps.
\item Ampleness at the limit. Stacks Project, Limits, Lemma limits-lemma-limit-ample. Gives \$\textbackslash{}mathcal\{L\}\_\{i, y\_i\}\$ ample for some \$i\$.
\item Ampleness in a neighbourhood, finite-presentation case. Stacks Project, Divisors on Spaces, Lemma spaces-divisors-lemma-ample-in-neighbourhood. Supplies the open \$V\_i\$ on the model.
\item Ampleness under affine pullback and base change. Stacks Project, Morphisms, Lemmas morphisms-lemma-ample-base-change and morphisms-lemma-pullback-ample-tensor-relatively-ample. Named by the author as the ingredients of the final pullback step.
\end{itemize}\medskip
\ifdefined\hypertarget\hypertarget{sourcepointer-1}{}\fi
\noindent\textit{Source reference.} Original reference 0D31: sources/stacks-chow/tags.txt, source lines 16401--16401. Cross-chapter standard prerequisite cited by the author inside the counted span; it is used as a pointer and its own proof is not part of this item.\par\medskip
\ifdefined\hypertarget\hypertarget{sourcepointer-2}{}\fi
\noindent\textit{Source reference.} Original reference 0D36: sources/stacks-chow/tags.txt, source lines 16406--16406. Cross-chapter standard prerequisite cited by the author inside the counted span; it is used as a pointer and its own proof is not part of this item.\par\medskip
\ifdefined\hypertarget\hypertarget{sourcepointer-3}{}\fi
\noindent\textit{Source reference.} Original lemma 0D3C: sources/stacks-chow/tags.txt, source lines 16412--16412. Cross-chapter standard prerequisite cited by the author inside the counted span; it is used as a pointer and its own proof is not part of this item.\par\medskip
\ifdefined\hypertarget\hypertarget{sourcepointer-4}{}\fi
\noindent\textit{Source reference.} Original reference 0A0X: sources/stacks-chow/tags.txt, source lines 12650--12650. Cross-chapter standard prerequisite cited by the author inside the counted span; it is used as a pointer and its own proof is not part of this item.\par\medskip
\ifdefined\hypertarget\hypertarget{sourcepointer-5}{}\fi
\noindent\textit{Source reference.} Original reference 07SQ: sources/stacks-chow/tags.txt, source lines 9910--9910. Cross-chapter standard prerequisite cited by the author inside the counted span; it is used as a pointer and its own proof is not part of this item.\par\medskip
\ifdefined\hypertarget\hypertarget{sourcepointer-6}{}\fi
\noindent\textit{Source reference.} Original reference 0D2X: sources/stacks-chow/tags.txt, source lines 16397--16397. Cross-chapter standard prerequisite cited by the author inside the counted span; it is used as a pointer and its own proof is not part of this item.\par\medskip
\ifdefined\hypertarget\hypertarget{sourcepointer-7}{}\fi
\noindent\textit{Source reference.} Original reference 07SR: sources/stacks-chow/tags.txt, source lines 9911--9911. Cross-chapter standard prerequisite cited by the author inside the counted span; it is used as a pointer and its own proof is not part of this item.\par\medskip
\ifdefined\hypertarget\hypertarget{sourcepointer-8}{}\fi
\noindent\textit{Source reference.} Original reference 09MT: sources/stacks-chow/tags.txt, source lines 12191--12191. Cross-chapter standard prerequisite cited by the author inside the counted span; it is used as a pointer and its own proof is not part of this item.\par\medskip
\ifdefined\hypertarget\hypertarget{sourcepointer-9}{}\fi
\noindent\textit{Source reference.} Original reference 0D3A: sources/stacks-chow/tags.txt, source lines 16410--16410. Cross-chapter standard prerequisite cited by the author inside the counted span; it is used as a pointer and its own proof is not part of this item.\par\medskip
\ifdefined\hypertarget\hypertarget{sourcepointer-10}{}\fi
\noindent\textit{Source reference.} Original reference 0893: sources/stacks-chow/tags.txt, source lines 10483--10483. Cross-chapter standard prerequisite cited by the author inside the counted span; it is used as a pointer and its own proof is not part of this item.\par\medskip
\ifdefined\hypertarget\hypertarget{sourcepointer-11}{}\fi
\noindent\textit{Source reference.} Original reference 0892: sources/stacks-chow/tags.txt, source lines 10482--10482. Cross-chapter standard prerequisite cited by the author inside the counted span; it is used as a pointer and its own proof is not part of this item.\par\medskip
\begin{lemma}
\label{lemma-ample-in-neighbourhood}
Let $S$ be a scheme.
Let $f : X \to Y$ be a proper morphism of algebraic spaces over $S$.
Let $\mathcal{L}$ be an invertible $\mathcal{O}_X$-module.
There exists an open subspace $V \subset Y$ characterized by
the following property:
A morphism $Y' \to Y$ of algebraic spaces factors
through $V$ if and only if the pullback $\mathcal{L}'$
of $\mathcal{L}$ to $X' = Y' \times_Y X$ is ample on $X'/Y'$
(as in Divisors on Spaces, Definition
\ref{spaces-divisors-definition-relatively-ample}).
\end{lemma}

\begin{proof}
Suppose that the lemma holds whenever $Y$ is a scheme.
Let $U$ be a scheme and let $U \to Y$ be a surjective \'etale morphism.
Let $R = U \times_Y U$ with projections $t, s : R \to U$.
Denote $X_U = U \times_Y X$ and $\mathcal{L}_U$ the pullback.
Then we get an open subscheme $V' \subset U$ as in the lemma
for $(X_U \to U, \mathcal{L}_U)$. By the functorial characterization
we see that $s^{-1}(V') = t^{-1}(V')$.
Thus there is an open subspace $V \subset Y$ such that
$V'$ is the inverse image of $V$ in $U$.
In particular $V' \to V$ is surjective \'etale and we
conclude that $\mathcal{L}_V$ is ample on $X_V/V$
(Divisors on Spaces, Lemma \ref{spaces-divisors-lemma-relatively-ample-local}).
Now, if $Y' \to Y$ is a morphism such that
$\mathcal{L}'$ is ample on $X'/Y'$, then
$U \times_Y Y' \to Y'$ must factor through $V'$
and we conclude that $Y' \to Y$ factors through $V$.
Hence $V \subset Y$ is as in the statement of the lemma.
In this way we reduce to the case dealt with in the next
paragraph.

\medskip\noindent
Assume $Y$ is a scheme. Since the question is local on $Y$
we may assume $Y$ is an affine scheme. We will show the
following:
\begin{enumerate}
\item[(A)] If $\Spec(k) \to Y$ is a morphism such that
$\mathcal{L}_k$ is ample on $X_k/k$, then there is an
open neighbourhood $V \subset Y$ of the image of $\Spec(k) \to Y$
such that $\mathcal{L}_V$ is ample on $X_V/V$.
\end{enumerate}
It is clear that (A) implies the truth of the lemma.

\medskip\noindent
Let $X \to Y$, $\mathcal{L}$, $\Spec(k) \to Y$ be as in (A).
By Lemma \ref{lemma-descending-property-ample}
we may assume that $k = \kappa(y)$ is the residue field of a point $y$
of $Y$.

\medskip\noindent
As $Y$ is affine we can find a directed set $I$ and an
inverse system of morphisms $X_i \to Y_i$ of algebraic spaces
with $Y_i$ of finite presentation over $\mathbf{Z}$, with affine
transition morphisms $X_i \to X_{i'}$ and $Y_i \to Y_{i'}$,
with $X_i \to Y_i$ proper and of finite presentation, and such that
$X \to Y = \lim (X_i \to Y_i)$. See Limits of Spaces, Lemma
\ref{spaces-limits-lemma-proper-limit-of-proper-finite-presentation-noetherian}.
After shrinking $I$ we may assume $Y_i$ is an (affine) scheme for all $i$,
see Limits of Spaces, Lemma \ref{spaces-limits-lemma-limit-is-affine}.
After shrinking $I$ we can assume we have a compatible system of
invertible $\mathcal{O}_{X_i}$-modules $\mathcal{L}_i$
pulling back to $\mathcal{L}$, see
Limits of Spaces, Lemma \ref{spaces-limits-lemma-descend-invertible-modules}.
Let $y_i \in Y_i$ be the image of $y$.
Then $\kappa(y) = \colim \kappa(y_i)$.
Hence $X_y = \lim X_{i, y_i}$ and after shrinking $I$ we
may assume $X_{i, y_i}$ is a scheme for all $i$, see
Limits of Spaces, Lemma \ref{spaces-limits-lemma-limit-is-scheme}.
Hence for some $i$ we have $\mathcal{L}_{i, y_i}$
is ample on $X_{i, y_i}$ by
Limits, Lemma \ref{limits-lemma-limit-ample}.
By Divisors on Spaces, Lemma \ref{spaces-divisors-lemma-ample-in-neighbourhood}
we find an open neighbourhood
$V_i \subset Y_i$ of $y_i$ such that
$\mathcal{L}_i$ restricted to $f_i^{-1}(V_i)$
is ample relative to $V_i$.
Letting $V \subset Y$ be the inverse image of
$V_i$ finishes the proof (hints: use
Morphisms, Lemma \ref{morphisms-lemma-ample-base-change} and
the fact that $X \to Y \times_{Y_i} X_i$ is affine
and the fact that the pullback of an
ample invertible sheaf by an affine morphism is ample by
Morphisms, Lemma \ref{morphisms-lemma-pullback-ample-tensor-relatively-ample}).
\end{proof}

\end{document}
